\section{Research Question, Hypothesis, and Study Design}
\label{sec:design}

\subsection{Research question and hypothesis}
\textbf{Research question.} For a facial expression classifier trained on simulated Emo-DVS events, does capping the influence of high-count temporal bins improve clip-level macro-F1 on real NEFER events without training or tuning on NEFER?

\textbf{Hypothesis.} Compared with the same classifier using unscaled event counts, adaptive scaling will improve mean NEFER macro-F1 by at least 3 percentage points while reducing mean Emo-DVS test macro-F1 by at most 2 points. These thresholds define practical success in advance. The prediction is that the spatial pattern, polarity, and sequence of facial changes transfer better than the absolute count of events per bin. The alternative is that high counts are useful expression cues, in which case the method may suppress genuine facial movement. A smaller gain, no gain, or a source penalty beyond the stated margin will not meet the hypothesis.

\subsection{Data, labels, and leakage controls}
\textbf{Simulated source.} Emo-DVS contains 9,024 simulated emotion clips and 4,042 real facial action-unit recordings \cite{chen2026}. We will use only compatible simulated emotion clips from the Ekman6* benchmark, after verifying the released labels and sequence format. The action-unit data are excluded from the main study because they use different labels and would introduce real sensor data into source training. The six-class vocabulary is anger, disgust, fear, happiness (including \emph{joy}), sadness, and surprise. Start from the authors' 70/10/20 train/validation/test protocol, which groups all clips derived from the same original video. Audit the released partitions for repeated originals and known identities. If they overlap or the manifest is absent, construct grouped splits before training, keeping the same original video and any known shared participant in a single partition. When subject IDs are unavailable, report the split as disjoint by original video and document the remaining risk of identity overlap.

\textbf{Real target.} NEFER provides paired RGB and event recordings of participants reacting to video stimuli \cite{berlincioni2023}. Use the six participants in the authors' validation partition as this study's locked target test set, and reserve participants in the official training partition solely for the optional real-data comparison. Use the participants' reported labels for the primary analysis; exclude \emph{None}, \emph{Contempt}, and all other labels outside the shared vocabulary under a rule fixed before scoring. Report the retained counts by class and participant, and use the expected labels on the same eligible clips as a separate diagnostic when both labels belong to the shared set. No NEFER data will be used to select the primary model, early-stopping epoch, scaling threshold, or evaluation rules. The goal is to predict dataset annotations from visible facial changes; the labels do not verify an individual's internal emotional state.

\textbf{Access and contingency.} The Emo-DVS authors describe release to verified researchers under a data-use agreement \cite{chen2026}; the study depends on authorized access to raw events or equivalent polarity-separated counts. The \href{https://github.com/miccunifi/NEFER}{NEFER repository} links to the recordings but does not state an explicit dataset license in its README. Confirm the applicable permissions for both datasets and record the versions, terms, and redistribution restrictions before using the recordings. If Emo-DVS access or the required event format is unavailable, and the NEFER terms permit it, convert paired NEFER RGB clips from the official training participants with v2e \cite{hu2021}. Use separate participant groups for simulated training, validation, and simulated testing within that partition; retain the same six held-out participants for the real test. Label this a \emph{within-NEFER} sim-to-real experiment, with no matching RGB or simulated version of any real test clip used in development.

\subsection{Baseline and one proposed change}
Use provided face regions where available, with a fixed rule that takes one crop per clip and avoids creating events by moving the crop frame by frame. If usable annotations are missing, use a fixed, label-blind crop policy for the affected dataset and log it before evaluation. Use only events as network input. Divide each clip into eight equal-duration bins across its own time range. Project events from the crop onto a $224\times224$ grid and count positive and negative events separately. For bin $b$, let $C_b^+$ and $C_b^-$ be the two count images and let $N_b=\sum_{x,y}(C_b^++C_b^-)$ be the total count. The conventional event-frame baseline stacks the positive count, negative count, and their sum as three channels, applies fixed log compression, and encodes each bin using a frozen DINOv2 ViT-S/14 backbone \cite{oquab2024}. A small gated recurrent unit reads the eight bin embeddings in order and a linear head predicts one of the six labels. Cache the frozen embeddings and train only the temporal head to keep the comparison feasible.

The proposed intervention is bounded per-bin count scaling. Before log compression, multiply both polarity channels in each bin by the same factor:
\begin{equation}
  s_b = \max\!\left(\tfrac{1}{4},
       \frac{R}{\max(R,N_b)}\right),\qquad
  \widetilde{C}_b^{\pm}=s_b C_b^{\pm},
  \label{eq:scale}
\end{equation}
where $R$ is the 75th percentile of nonempty bin counts on the simulated training split. Recompute the summed channel from the scaled counts. The transform leaves sparse and empty bins unchanged, preserves event locations and polarity balance before compression, and limits the attenuation of dense bins to a factor of four. It is intended to reduce sensitivity to broad changes in event density while retaining the order of motion across bins. It cannot distinguish dense expression motion from illumination bursts or remove sensor noise in individual pixels. The backbone weights, binning, preprocessing outside Eq.~\eqref{eq:scale}, optimizer, class weights, and training budget will be held fixed across the main comparison.

\subsection{Controlled comparisons and evaluation}
Train the baseline and the proposed model with the same three random seeds and select checkpoints by simulated validation macro-F1. Score each locked model on both simulated test clips and real test clips. The simulated scores measure in-domain performance; the paired real scores give the primary estimate of the effect of scaling. Also report each model's simulated-minus-real score difference as a descriptive measure of transfer. For the ablation, replace $N_b$ in Eq.~\eqref{eq:scale} with the mean count across the eight bins of the clip; using one scale for the entire clip tests whether any benefit comes from adapting within a clip or from scaling generally.

If time and class coverage allow, take the locked source baseline and update only its linear head using clips from five preselected participants in NEFER's official training partition. Keep the same target test set; report this as a separate limited-supervision comparison. If a matching evTransFER implementation is unavailable under the same source data, label set, and split, use the published results to discuss the method rather than as an experimental ranking. The primary comparison needs one main baseline and one change to its representation.

Use clip-level macro-F1 as the primary metric so every class has equal weight, with accuracy, per-class precision and recall, and row-normalized confusion matrices as secondary outputs. Report results per seed and their mean. Assess whether any single NEFER participant drives the gain by recomputing the paired macro-F1 difference while omitting one participant at a time, and state when an omission leaves a class unrepresented. Inspect errors by event density; if time permits, stress-test both models on the same fixed random event-drop and event-injection perturbations. These perturbations test sensitivity to the chosen corruptions and do not replace measured illumination trials. Freeze split manifests, class mappings, settings, and evaluation code before testing; save aggregate scores, predictions where allowed, and model configurations. The six-participant target split limits how broadly either a gain or a negative result can be interpreted.
